% sec:adm-foliation — Foliation of Spacetime
\section{Foliation of Spacetime}
\label{sec:adm-foliation}

Einstein's field equations (\cref{sec:gr-efe}) describe spacetime as a
single four-dimensional entity --- they are covariant, making no
distinction between space and time.  A numerical solver, by contrast,
needs initial data on a spatial surface and a rule for marching that data
forward.  The bridge between the two viewpoints is the \emph{foliation}: a
decomposition of the four-dimensional manifold into a one-parameter family
of three-dimensional spatial slices, each representing the state of the
gravitational field at one instant of a chosen time parameter.  Three
objects control the decomposition: a time function whose level sets define
the slices, a unit normal vector that threads them together, and an
induced metric that measures distances within each slice.
\histnote{The 3+1 decomposition was developed by Arnowitt, Deser, and
Misner \cite{ADM1962} in the context of canonical quantum gravity.  Its
reformulation as an initial-value problem for numerical work was
crystallised by York \cite{York1979}.}

\paragraph{The time function.}

The construction begins with a smooth scalar field $t$ on the spacetime
manifold $\mathcal{M}$.  Each level set
$\Sigma_t = \{p \in \mathcal{M} : t(p) = \text{const}\}$ is a
three-dimensional hypersurface.  For these slices to represent genuine
spatial surfaces, the gradient must be timelike everywhere:
%
\begin{equation}
  \inversemetric \,(\cd{\mu} t)\,(\cd{\nu} t) < 0 \,.
  \label{eq:adm-time-function}
\end{equation}
%
When this condition holds, the normal to each slice is timelike, so by the
signature of the Lorentzian metric its orthogonal complement is spacelike
--- every vector tangent to $\Sigma_t$ has positive norm, making
$\Sigma_t$ a suitable surface on which to pose initial data.  The family
$\{\Sigma_t\}$ fills the spacetime with
non-intersecting leaves: every event lies on exactly one slice, and the
parameter $t$ labels which one.
\warnnote{Not every spacetime admits a global foliation.  The existence of
a smooth time function satisfying \cref{eq:adm-time-function} requires
global hyperbolicity --- a condition met by every physically reasonable
spacetime used in numerical relativity, but violated by, for example,
the maximal Kerr extension through the Cauchy horizon
(\cref{sec:kerr-horizons}).}

\paragraph{The unit normal.}

The gradient $\cd{\mu} t$ is perpendicular to each slice $\Sigma_t$ by
construction --- any vector $V^\mu$ tangent to the slice satisfies
$V^\mu\,\cd{\mu} t = 0$ because $t$ is constant along $\Sigma_t$.  But
the gradient need not have unit norm.  Normalising it produces the unit
normal covector
%
\begin{equation}
  n_\mu = -\lapse\,\cd{\mu} t \,,
  \label{eq:adm-normal-covector}
\end{equation}
%
where the positive scalar
%
\begin{equation}
  \lapse \equiv \bigl(-\inversemetric\,(\cd{\mu} t)\,(\cd{\nu} t)
    \bigr)^{-1/2}
  \label{eq:adm-lapse-normalisation}
\end{equation}
%
ensures the normalisation $n_\mu\,n^\mu = -1$.  Substituting
\cref{eq:adm-normal-covector} into $n_\mu\,n^\mu$ gives
$n_\mu\,n^\mu = \lapse^2\,\inversemetric\,(\cd{\mu} t)(\cd{\nu} t)
= \lapse^2 \cdot (-\lapse^{-2}) = -1$ --- the normalisation is built in
by design.  With the $(-,+,+,+)$ signature, raising the index on the
timelike covector $\cd{\mu} t$ reverses its time orientation; the explicit
minus sign in \cref{eq:adm-normal-covector} compensates, making $n^\mu$
future-directed.

The scalar $\lapse$ has a direct physical meaning.  An observer who moves
along $n^\mu$ --- perpendicular to each slice, with zero spatial velocity
--- ages by
%
\begin{equation}
  \d\tau = \lapse\,\d t
  \label{eq:adm-proper-time}
\end{equation}
%
between neighbouring slices.  The quantity $\lapse$ is the \emph{lapse
function}: it ``lapses'' proper time between slices, translating the
coordinate label $t$ into the physical clock reading $\tau$.
\cref{sec:adm-lapse-shift} develops $\lapse$ as one of two gauge variables
that encode the freedom in threading coordinates through the foliation.
\physnote{The observer who moves along $n^\mu$ is the \emph{Eulerian
observer} of numerical relativity --- the fiducial reference frame for
interpreting ADM quantities.  The extrinsic curvature, for instance,
measures how the spatial slice curves as seen by this observer.}

\paragraph{The spatial projection.}

To do geometry within a slice we need a systematic way to extract the
spatial part of any spacetime quantity --- discarding its component along
the normal while preserving the rest.  The \emph{spatial projection
operator} performs this split:
%
\begin{equation}
  \gamma^\mu{}_\nu = \delta^\mu{}_\nu + n^\mu\,n_\nu \,.
  \label{eq:adm-projector}
\end{equation}
%
Applied to any vector, $\gamma^\mu{}_\nu\,V^\nu$ returns the part tangent
to $\Sigma_t$ and annihilates the component along $n^\mu$.  Two properties
confirm that this is a genuine projection.  First, it is orthogonal to the
normal:
%
\begin{equation}
  \gamma^\mu{}_\nu\,n^\nu
  = n^\mu + n^\mu\,\underbrace{n_\nu\,n^\nu}_{-1} = 0 \,.
  \label{eq:adm-projector-orthogonal}
\end{equation}
%
Second, it is idempotent --- projecting twice gives the same result as
projecting once:
%
\begin{equation}
  \gamma^\mu{}_\alpha\,\gamma^\alpha{}_\nu
  = \delta^\mu{}_\nu + 2\,n^\mu\,n_\nu
    + n^\mu\,\underbrace{n_\alpha\,n^\alpha}_{-1}\,n_\nu
  = \gamma^\mu{}_\nu \,.
  \label{eq:adm-projector-idempotent}
\end{equation}
%
Together these properties guarantee that $\gamma^\mu{}_\nu$ cleanly
separates every vector into a spatial part and a normal part with no
residual mixing: the spatial component lives entirely within $\Sigma_t$,
and projecting it again changes nothing.

\begin{figure}[htbp]
  \centering
  % adm-foliation.tex — Foliation of spacetime into spacelike hypersurfaces
% Inputable TikZ fragment for fig:adm-foliation
% Requires: tikz with calc, arrows.meta (loaded by ehbook.cls)
%
% Schematic spacetime diagram showing two neighbouring spatial slices
% Sigma_t and Sigma_{t+dt}, the unit normal n^mu, the proper-time gap
% alpha dt, and a tangent vector V^mu.  The gap between slices varies
% across the surface to show that alpha is spatially non-uniform.
%
% Corners of each sheet (perspective parallelogram):
%   Front-to-back perspective offset: dx = -0.7, dy = +1.5
%   Lower: FL=(0.5,0) FR=(7.0,0.20) BR=(6.3,1.70) BL=(-0.2,1.50)
%   Upper: FL=(0.5,2.50) FR=(7.0,3.40) BR=(6.3,4.90) BL=(-0.2,4.00)
%   Gap ranges from ~2.5 (left) to ~3.2 (right).
%
\begin{tikzpicture}[
    >={Stealth[length=5pt]},
    font=\footnotesize,
  ]

  % ── Lower sheet (Σ_t) ─────────────────────────────────────────────

  % Fill
  \fill[EHMarginGray, opacity=0.06]
    plot[smooth, tension=0.6] coordinates {
      (0.5, 0.0) (2.5, -0.08) (4.5, 0.06) (7.0, 0.20)}
    -- (6.3, 1.70) -- (-0.2, 1.50) -- cycle;

  % Front edge (closest to viewer)
  \draw[EHMarginGray, semithick]
    plot[smooth, tension=0.6] coordinates {
      (0.5, 0.0) (2.5, -0.08) (4.5, 0.06) (7.0, 0.20)};

  % Back edge
  \draw[EHMarginGray, thin, opacity=0.4]
    (-0.2, 1.50) -- (6.3, 1.70);

  % Side edges
  \draw[EHMarginGray, thin, opacity=0.4]
    (0.5, 0.0) -- (-0.2, 1.50)  (7.0, 0.20) -- (6.3, 1.70);

  % Coordinate gridlines — transverse (depth 1/3, 2/3)
  \draw[EHMarginGray, very thin, opacity=0.15]
    (0.27, 0.50) -- (6.77, 0.70)
    (0.03, 1.00) -- (6.53, 1.20);

  % Coordinate gridlines — longitudinal (width 1/4, 1/2, 3/4)
  \draw[EHMarginGray, very thin, opacity=0.15]
    (2.13, 0.05) -- (1.43, 1.55)
    (3.75, 0.10) -- (3.05, 1.60)
    (5.38, 0.15) -- (4.68, 1.65);

  % Label
  \node[font=\small, text=EHMarginGray, anchor=west]
    at (7.15, 0.20) {$\Sigma_t$};


  % ── Upper sheet (Σ_{t+dt}) ────────────────────────────────────────

  \fill[EHMarginGray, opacity=0.06]
    plot[smooth, tension=0.6] coordinates {
      (0.5, 2.50) (2.5, 2.44) (4.5, 2.96) (7.0, 3.40)}
    -- (6.3, 4.90) -- (-0.2, 4.00) -- cycle;

  \draw[EHMarginGray, semithick]
    plot[smooth, tension=0.6] coordinates {
      (0.5, 2.50) (2.5, 2.44) (4.5, 2.96) (7.0, 3.40)};

  \draw[EHMarginGray, thin, opacity=0.4]
    (-0.2, 4.00) -- (6.3, 4.90);

  \draw[EHMarginGray, thin, opacity=0.4]
    (0.5, 2.50) -- (-0.2, 4.00)  (7.0, 3.40) -- (6.3, 4.90);

  % Gridlines — transverse
  \draw[EHMarginGray, very thin, opacity=0.15]
    (0.27, 3.00) -- (6.77, 3.90)
    (0.03, 3.50) -- (6.53, 4.40);

  % Gridlines — longitudinal
  \draw[EHMarginGray, very thin, opacity=0.15]
    (2.13, 2.73) -- (1.43, 4.23)
    (3.75, 2.95) -- (3.05, 4.45)
    (5.38, 3.18) -- (4.68, 4.68);

  \node[font=\small, text=EHMarginGray, anchor=west]
    at (7.15, 3.40) {$\Sigma_{t+\mathrm{d}t}$};


  % ── Point P on lower sheet ────────────────────────────────────────
  \coordinate (P) at (3.3, 0.60);
  \fill (P) circle (1.5pt);
  \node[font=\scriptsize, anchor=north east, inner sep=2pt]
    at (P) {$P$};


  % ── Unit normal n^μ (EHJetBlue) ───────────────────────────────────
  %
  % Vertical from P to upper sheet.  At x = 3.3, the upper-sheet
  % surface (same depth fraction) sits at approximately y ≈ 3.38.
  \draw[->, EHJetBlue, very thick, shorten >=2pt]
    (P) -- (3.3, 3.38);
  \node[font=\scriptsize, text=EHJetBlue, anchor=east, inner sep=2pt]
    at (3.25, 3.15) {$n^\mu$};


  % ── Proper-time gap α dt (EHAccretionGold) ────────────────────────
  %
  % Measurement bracket to the right of the normal arrow.
  \draw[EHAccretionGold, thin] (3.50, 0.60) -- (3.70, 0.60);
  \draw[EHAccretionGold, thin] (3.50, 3.38) -- (3.70, 3.38);
  \draw[EHAccretionGold, thin, <->] (3.60, 0.60) -- (3.60, 3.38);
  \node[font=\scriptsize, text=EHAccretionGold, anchor=west, inner sep=1pt]
    at (3.75, 1.99) {$\alpha\,\mathrm{d}t$};


  % ── Tangent vector V^μ (EHAccretionGold) ──────────────────────────
  \draw[->, EHAccretionGold, very thick]
    (P) -- ++(1.5, 0.06);
  \node[font=\scriptsize, text=EHAccretionGold, anchor=south west, inner sep=1pt]
    at (4.65, 0.68) {$V^\mu$};


  % ── Right-angle marker at P ───────────────────────────────────────
  %
  % Small square between n^μ (vertical) and V^μ (horizontal).
  \draw[thin] (3.3, 0.72) -- (3.42, 0.72) -- (3.42, 0.60);

\end{tikzpicture}

  \caption[Foliation of spacetime into spacelike hypersurfaces]{%
    Foliation of spacetime into spacelike hypersurfaces.
    Two neighbouring slices $\Sigma_t$ and $\Sigma_{t+\d t}$ are
    separated by a proper-time interval $\alpha\,\d t$ along the unit
    normal $n^\mu$.  A tangent vector $V^\mu$ lies within $\Sigma_t$,
    satisfying $n_\mu V^\mu = 0$.  The lapse function $\alpha$ varies
    across the slice, so the proper-time gap between slices is not
    uniform --- the geometry stretches time differently at different
    spatial locations.}
  \label{fig:adm-foliation}
\end{figure}

\paragraph{The induced metric.}

The fully covariant form of the projector serves double duty as the
\emph{induced metric} (or first fundamental form) on $\Sigma_t$:
%
\begin{equation}
  \gamma_{\mu\nu} = \metric + n_\mu\,n_\nu \,.
  \label{eq:adm-induced-metric}
\end{equation}
%
For two vectors $U^\mu$ and $V^\mu$ both tangent to the slice,
$n_\mu\,U^\mu = 0$ and $n_\nu\,V^\nu = 0$, so
$\gamma_{\mu\nu}\,U^\mu\,V^\nu = \metric\,U^\mu\,V^\nu$ --- the induced
metric agrees with the full spacetime metric on spatial vectors.
In coordinates $(t, x^i)$ adapted to the foliation, the spatial components
$\gamma_{ij}$ form a non-degenerate three-metric $\spatialmetric$ that
carries
all the information about the slice's intrinsic geometry: its own
Christoffel symbols, Riemann tensor, and curvature invariants are built
from $\gamma_{ij}$ and its spatial derivatives exactly as the spacetime
quantities are built from $\metric$ (\cref{sec:dg-curvature}).
\xrefnote{The induced metric captures the slice's \emph{intrinsic}
geometry --- what a being confined to $\Sigma_t$ could measure.  The
\emph{extrinsic} curvature $\extrinsic$
(\cref{sec:adm-extrinsic-curvature}) captures how the slice bends within
the ambient spacetime.  Together, $\spatialmetric$ and $\extrinsic$
constitute the initial data for the ADM evolution.}

The foliation, the unit normal, and the induced metric are the geometric
foundation for the 3+1 decomposition.  The spacetime metric $\metric$ has
been split into $\gamma_{\mu\nu}$ (geometry within each slice) and $n_\mu$
(geometry between slices).  What remains is to specify how coordinates are
laid across this structure --- how much proper time elapses between slices
(the lapse $\lapse$) and how the spatial coordinates shift from one slice
to the next (the shift $\shift$).  These are the gauge variables of the
next section.
